\documentclass[letterpaper]{article} % DO NOT CHANGE THIS
\usepackage[submission]{aaai2027}  % DO NOT CHANGE THIS
\usepackage[hyphens]{url}  % DO NOT CHANGE THIS
\usepackage{graphicx} % DO NOT CHANGE THIS
\urlstyle{rm} % DO NOT CHANGE THIS
\def\UrlFont{\rm}  % DO NOT CHANGE THIS
\usepackage{natbib}  % DO NOT CHANGE THIS AND DO NOT ADD ANY OPTIONS TO IT
\usepackage{caption} % DO NOT CHANGE THIS AND DO NOT ADD ANY OPTIONS TO IT
\frenchspacing  % DO NOT CHANGE THIS
\usepackage{booktabs}
\usepackage{amsmath}
\usepackage{mdframed}

\pdfinfo{
/TemplateVersion (2027.1)
}

\setcounter{secnumdepth}{1} % Enabled (AAAI-permitted) so in-text "Appendix A"/"Appendix B" references resolve to real section letters after \appendix.

\newmdenv[
  linewidth=0.6pt, linecolor=black!55, backgroundcolor=black!3,
  innertopmargin=6pt, innerbottommargin=6pt, innerleftmargin=8pt, innerrightmargin=8pt,
  skipabove=8pt, skipbelow=4pt
]{harnessbox}

\title{EdgeCase: An Assurance Harness for Conflicting Governance Risks in Agentic Systems}
\author{
    Anonymous Submission
}
\affiliations{
}

\begin{document}

\maketitle

\begin{abstract}
Agentic AI systems must satisfy security, privacy, fairness, accessibility, care, and efficiency obligations at once, yet assurance methods still evaluate them independently -- implicitly treating a win on one metric as a win for the whole system. We show this assumption fails and introduce the \emph{assurance harness}: a bounded, auditable routing layer that detects when governance obligations collide and selects a mitigation, decoupled from whatever detector drives it -- the way a test harness decouples test logic from the system under test. We instantiate this as EdgeCase and establish three findings on a 1,260-instance, seven-domain benchmark: ten regenerated seeds, cross-model evaluation, external validity against two independent agent-safety benchmarks and the largest published AI risk taxonomy, and a blinded human-annotator study. First, governance policies that commit to one fixed action are structurally blind to most collision types they encounter, independent of how externalities are scored. Second, which specific policy has the lowest aggregate externality is not assumption-free: reasonable alternative scoring and weighting choices reverse specific pairwise rankings. Third, the harness's reliability is separable from its detector: a correctly-specified deterministic detector matches a hindsight-optimal oracle and beats a capable off-the-shelf LLM classifier on the identical routing table -- but that correctness is conditional on a specification-completeness burden that is itself error-prone: two domains' trigger vocabularies silently collapsed detection to 36--48\% before we caught it. Assurance for agentic AI requires an auditable harness that actively manages competing obligations, not a single optimized target.
\end{abstract}

\begin{links}
    \link{Code and reproducibility artifacts}{https://github.com/anon0404/edgecase}
\end{links}

\section{Introduction}
Large language models (LLMs) are increasingly embedded within agentic systems that perform multi-step reasoning, memory management, retrieval, planning, escalation, delegation, and tool-mediated execution \citep{yao2023react,schick2023toolformer,park2023generative,sapkota2025aiagents,oecd2026landscape}. These systems are now deployed in high-stakes settings including healthcare triage, financial services, crisis assistance, education, enterprise operations, and public-sector workflows \citep{hughes2025aiagents,brohi2025landscape}. Unlike single-turn conversational systems, agentic architectures operate through persistent orchestration pipelines, where governance decisions emerge sequentially across interacting components \citep{yao2023react,wu2024autogen,sapkota2025aiagents}.

Existing approaches to AI safety and governance typically evaluate security, fairness, privacy, accessibility, transparency, and efficiency risks independently of each other \citep{bai2022constitutional,ouyang2022training,mitchell2019modelcards,eu2024ai}. Security benchmarks evaluate jailbreak resistance and harmful-output suppression \citep{ganguli2022redteaming,anthropic2025rsp}; fairness evaluations measure disparate performance across demographic groups \citep{barocas2016disparate,selbst2019fairness}; privacy frameworks assess information leakage and retention risk \citep{dwork2014algorithmic}; operational metrics measure latency, throughput, and cost \citep{hughes2025aiagents}. Real-world agentic deployments rarely encounter these obligations in isolation, however. They increasingly operate under conditions where several governance requirements become active simultaneously and recommend incompatible actions \citep{syros2025saga,khan2025agentsafe}.

A crisis-support assistant may encounter language that resembles adversarial prompt injection while simultaneously containing a credible self-harm disclosure; a banking assistant might need stronger identity verification while serving a user facing linguistic or accessibility barriers; a healthcare workflow may improve oversight reliability through deeper review while substantially increasing latency, energy consumption, and cost \citep{eu2024ai,syros2025saga,huang2025aagate}. In each case, no single fixed policy is right for the case beside it.

These tensions expose a limitation in current AI assurance methodology, not just in current deployments. Safety interventions typically assume governance objectives can be independently optimized without inducing secondary harms elsewhere in the system \citep{bai2022constitutional,ouyang2022training}. In practice, strengthening one governance property often degrades another. We argue these failures are not isolated edge cases but structural conflicts between interacting governance obligations \citep{syros2025saga,huang2025aagate}: a localized safety improvement can look like progress on its own metric while simply redistributing governance burden to a domain nobody was measuring.

We call the space of these interacting obligations the \emph{governance tradeoff frontier}: rather than evaluating a system against one target metric, we model mitigation policies within a multi-objective space spanning security robustness, care preservation, accessibility burden, privacy exposure, oversight reliability, and energy consumption \citep{deb2001multiobjective}. Managing that frontier at runtime requires a component that can (i) detect when obligations collide and (ii) select a bounded, defensible response -- without hard-coding that response to any one detection mechanism, since a rule-based detector, a learned classifier, or a future replacement should all be able to plug into the same decision logic and produce an auditable outcome. We call this component an \emph{assurance harness}, by analogy to a software test harness: a test harness fixes the contract between a test suite and the system under test so the suite keeps working when the underlying implementation changes; an assurance harness fixes the contract between a bounded, auditable mitigation-routing table and whatever detector feeds it, so the routing logic -- the auditable part -- stays constant while the detector implementing it can be swapped, inspected, or upgraded independently. This decoupling is not just conceptually appealing: we demonstrate it directly in Section~\ref{sec:llm-detected} by swapping EdgeCase's detector for an off-the-shelf LLM without changing a single entry in its routing table.

This paper introduces EdgeCase, an implementation of an assurance harness for agentic AI. It instruments workflow traces, activates governance obligations from runtime signals, detects boundary collisions between obligations that recommend incompatible actions, and routes each collision through a bounded, auditable mitigation table instead of collapsing to binary allow-or-block behavior.

Our work is motivated by three research questions: (1) Do governance conflicts emerge systematically across agentic workflows and model ecosystems? (2) To what extent do single-objective mitigation strategies externalize harm into other governance domains? (3) Can conflict-aware routing reduce aggregate externalities relative to rigid safety policies, and does its reliability depend on the detector driving it?

The contributions of this paper are threefold. First, we show that governance policies committing to one fixed action are structurally blind to most conflict types they encounter -- an assumption-free result, independent of how externalities are scored -- while showing directly, not just acknowledging, that claims about which specific policy has the lowest aggregate cost are comparatively fragile and reverse under several independently-reasonable modeling choices. Second, we release EdgeCase-Benchmark-v1, a reusable methodology for stress-testing governance claims in agentic systems beyond this paper's specific results. Third, we show that an assurance harness's reusable artifact is its bounded, auditable mitigation-routing table, not the specific detector feeding it: a correctly-specified deterministic implementation of that table is more reliable than a capable off-the-shelf LLM classifier on the identical task -- not because LLM judgment is untrustworthy in absolute terms, but because a correctly-specified deterministic lookup need not fail at all, where the LLM's otherwise-small error rate remains nonzero. We report this alongside its cost: achieving that correctness requires verifying that an obligation registry's trigger vocabulary is complete across every domain, itself a nontrivial and error-prone specification burden, which we illustrate with the two domains (of seven) whose initial trigger vocabularies had gaps large enough to collapse detection to 36--48\% there.

\section{Related Work}
\paragraph{AI Safety and Alignment.} Contemporary AI safety research has focused on reducing harmful outputs, improving instruction following, increasing robustness to adversarial prompting, and aligning model behavior with human preferences or constitutional principles \citep{ouyang2022training,bai2022constitutional,ganguli2022redteaming}. Model documentation frameworks, algorithmic auditing, and governance standards operationalize these principles through reporting, risk-management processes, evaluation pipelines, and institutional accountability \citep{mitchell2019modelcards,gebru2021datasheets,raji2020closing,nist2023airmf,eu2024ai}. Recent agent-focused safety benchmarks, including AgentHarm and Agent-SafetyBench, show that LLM agents introduce safety risks beyond chatbot settings because agents execute multi-stage tasks, retain capabilities after jailbreaks, and act through external tools \citep{andriushchenko2024agentharm,zhang2024agentsafetybench}. We use both benchmarks directly as external validity checks in this paper (Appendix~\ref{app:external-validity}).

\paragraph{Responsible AI Governance.} Responsible AI research emphasizes fairness, accountability, transparency, privacy, explainability, contestability, and human oversight as central governance principles \citep{barocas2016disparate,mitchell2019modelcards,dwork2014algorithmic,selbst2019fairness,raji2020closing,nist2023airmf}. Existing frameworks, however, typically operationalize these concerns as separable dimensions evaluated through domain-specific rubrics. This is not a framing choice we impute to the field: the MIT AI Risk Repository \citep{slattery2026airisk}, a meta-review synthesizing 1,700+ risks from 74 frameworks into the most comprehensive systematic AI risk taxonomy publicly available, states directly that its taxonomies do not categorize risks by impact or likelihood and do not discuss interaction between risks -- naming risk-interaction modeling an explicit gap for future work rather than something the repository attempts. EdgeCase's harness treats governance as obligations that can be jointly active and mutually incompatible, rather than independently-tracked risk categories, which addresses precisely this gap (Appendix~\ref{app:taxonomy-validation} reports a bounded validation of EdgeCase's own registry against this repository and 16 real incident reports). Prior work shows governance interventions interacting in complex, sometimes unintended ways: fairness interventions can weaken privacy guarantees by requiring collection of sensitive attributes for bias assessment \citep{jayaraman2022revisiting,selbst2019fairness}; transparency mechanisms can improve accountability while increasing exploitability of enforcement logic \citep{kroll2017accountable,carlini2023extracting}; privacy-preserving approaches can remove contextual information needed for safeguarding or abuse detection \citep{dwork2014algorithmic,whittaker2021deep}; human oversight can introduce latency or uneven accessibility burden \citep{green2019disparate}. Three recent systems address adjacent governance problems: AGENTSAFE proposes a lifecycle governance framework \citep{khan2025agentsafe}, SAGA proposes cryptographic agent-to-agent access control \citep{syros2025saga}, and AAGATE proposes a NIST AI RMF-aligned control-plane platform \citep{huang2025aagate}. None implement a collision-and-mitigation harness in EdgeCase's sense -- their decision rules answer a different question (identity/access, not which of two active obligations should govern a case) -- so we compare EdgeCase against all three structurally rather than empirically in Appendix~\ref{app:baseline-comparison}.

\paragraph{Tradeoffs-aware AI Assurance.} Multi-objective optimization studies how systems search for Pareto-efficient solutions when objectives conflict \citep{deb2001multiobjective,miettinen1999nonlinear}. This literature is useful but insufficient for AI governance: many governance requirements are contextual obligations that remain normatively incompatible even inside a shared optimization space -- e.g., a privacy obligation to minimize retention can directly conflict with a safeguarding obligation to preserve evidence for escalation, where neither obligation is simply outweighed by the other. Responsible-ML work on fairness, accuracy, privacy, transparency, and robustness tradeoffs \citep{barocas2016disparate,selbst2019fairness,dwork2014algorithmic,jayaraman2022revisiting} largely evaluates them at the level of model outputs or decision rules, not how they emerge dynamically inside agentic workflows involving memory, retrieval, routing, tool use, and escalation. Recent agent benchmarks and orchestration frameworks show that consequential failures often emerge through sequential interactions, tool calls, and environment feedback rather than a single unsafe response \citep{wu2024autogen,yao2024taubench,yang2024sweagent,andriushchenko2024agentharm,zhang2024agentsafetybench}, motivating the runtime instrumentation and trace-based evaluation this paper adopts.

\section{Methodology}
\label{sec:methodology}
\begin{harnessbox}
\textbf{Definition 1 (Assurance Harness).} An assurance harness is a triple $H = (R, D, M)$: an obligation registry $R$ specifying which runtime signals trigger which governance obligations; a detector $D$ that maps observed signals to an activated-collision type from a closed set; and a mitigation router $M$ -- a bounded, auditable lookup from collision type to mitigation action, independent of $D$'s internal implementation. Only $D$'s I/O contract is fixed (signals in, one closed-set collision label out); $D$ itself is swappable. $R$ and $M$ are the harness's reusable, inspectable artifacts. EdgeCase (Sections~\ref{sec:results}--\ref{sec:llm-detected}) instantiates $H$ with a rule-based $D$ by default and demonstrates a second, LLM-based $D$ driving the identical $R, M$.
\end{harnessbox}

\paragraph{Governance Obligations.} We model governance as a dynamic system of interacting runtime obligations that activate conditionally throughout workflow execution. An obligation is represented as $O_i = (T_i, A_i, D_i, W_i)$, where $T_i$ denotes triggering conditions, $A_i$ a recommended governance action, $D_i$ the governance domain, and $W_i$ an obligation weight. Triggering conditions are evaluated over workflow state $S_t = (P_t, M_t, R_t, U_t, E_t)$ (prompt/conversational context, memory state, retrieval outputs, tool/routing state, and execution metadata, respectively): $T_i: S_t \rightarrow \{0,1\}$.

\paragraph{Boundary Collisions.} A boundary collision occurs when multiple activated obligations recommend incompatible governance actions. Letting $O_t = \{O_1, \ldots, O_n\}$ denote activated obligations at execution step $t$, a collision between two obligations is detected whenever their recommended actions form a registered incompatible-action pair. For example, $O_{\text{security}} = (T_{\text{jailbreak}}, \text{block}, \text{security})$ and $O_{\text{care}} = (T_{\text{self\_harm}}, \text{escalate}, \text{care})$ can activate together, producing a collision between \texttt{block} and \texttt{escalate}: the same signal set can activate multiple legitimate obligations recommending incompatible actions. The implementation registers seven such incompatible-action pairs, spanning thirteen obligations across ten governance domains (security, care, accessibility, privacy, safeguarding, environment, safety, memory, transparency, and fairness).

Because detection is an exact-match check against each obligation's registered trigger vocabulary, an obligation can only fire on signals explicitly listed as its triggers -- the registry's completeness is therefore part of the harness's specification, not an implementation detail. We treat the trigger vocabulary as a specification that must be verified, not merely authored: every case in a released benchmark is checked by running its signal set through the same detection function used at evaluation time before the case is committed, confirming its intended collision is actually reachable. Section~\ref{sec:results} and Appendix~\ref{app:taxonomy-validation} report what this verification step catches, and what it does not.

\paragraph{Governance Tradeoff Frontier and Aggregate Externality.} For workflow execution $k$, governance performance is represented as $G_k = (s_k, c_k, a_k, p_k, e_k)$ (security robustness, care preservation, accessibility burden, privacy exposure, and energy cost, respectively). We define aggregate governance externality as $X_k = \sum_i \lambda_i \Delta_i$, where $\Delta_i$ is the mitigation-induced burden on dimension $i$ and $\lambda_i$ a governance weighting coefficient. We report $X_k$ under a severity-ordered weighting ($\lambda = 0.30, 0.15, 0.20, 0.25, 0.10$ for care suppression, accessibility burden, privacy exposure, security risk, and energy cost, reflecting that care suppression risks physical or psychological harm to a vulnerable user while energy cost is an operational externality on the system operator rather than a harm to any stakeholder) and test every pairwise comparison's sensitivity to this and other modeling choices directly, in Section~\ref{sec:results} and Appendix~\ref{app:baseline-scoring-robustness}, rather than reporting one weighting as though it were the only defensible choice.

\paragraph{Conflict-aware Mitigation Routing.} EdgeCase's router $M$ selects a bounded mitigation conditioned on the specific obligation pair in tension for a case -- e.g., \texttt{constrain\_and\_escalate} for block-versus-escalate, \texttt{adaptive\_verification} for verify-versus-accessibility, \texttt{typed\_memory} for memory-care-versus-memory-poisoning -- rather than a binary allow-or-block decision. Four rigid single-action baseline policies (Strict Block, Always Escalate, Always Verify, Maximum Review) are compared against this adaptive router; each commits to one fixed action regardless of case content, i.e., each policy operates without a detector or a router at all.

\section{Experimental Design}
\paragraph{Benchmark Construction.} EdgeCase-Benchmark-v1 comprises 1,260 synthetic paired governance-conflict instances across seven deployment domains: crisis-support, banking/fintech, healthcare triage, enterprise copilot, education, moderation, and public-benefits. Each domain contributes exactly 180 instances (14.3\% of the benchmark) by deliberate design: equal domain representation lets policy comparisons be read without conflating a policy's performance with how often its worst-case domain happens to appear in the sample. We do not claim this reflects real-world deployment prevalence, which would require production traffic data outside this benchmark's scope. Each instance pairs a domain-specific prompt template with signals drawn from two competing obligation triggers, an expected collision type, and a reference bounded mitigation; every instance's signal set is verified against the detector at generation time (Methodology), so its intended collision type is guaranteed reachable. The benchmark provides standardized 70/15/15 train/validation/test splits. Domain scenarios are grounded in prior governance literature: crisis-support in mitigation failures in crisis-support LLM deployment \citep{miner2020suicide}; banking/fintech in identity-verification and financial-inclusion tensions \citep{fsb2017ai,bis2024ai}; healthcare triage in human-oversight and reliability requirements for medical AI \citep{topol2019highperformance,singhal2025llms,hoffmann2024llms}; enterprise copilot in memory-safety risk in agentic deployments \citep{openai2024memory,jagielski2023poisoning}; education in algorithmic bias and adaptive-learning equity \citep{holmes2022ethics,baker2019fairness}; moderation in transparency-versus-adversarial-adaptation tensions \citep{gillespie2018custodians,kolt2025governing}; and public-benefits in data-minimization and safeguarding tensions in digital public services \citep{ec2023digital,oecd2024safeguarding,oecd2024classification}. Broader deployment context for agentic systems is surveyed elsewhere \citep{staufer2026index,yuan2024benchmarking,xu2024bagoftricks,wang2025agentspec,whittaker2025agentic,doshivelez2017rigorous}.

\paragraph{Model-backed Evaluation.} To test whether governance conflicts generalize across heterogeneous model ecosystems, we run experiments using three model environments: Claude Sonnet, Gemini 2.5 Pro, and a locally-hosted Qwen2.5-7B via Ollama. Each provider is evaluated on a stratified 140-case sample (20 cases per domain), all executed through a shared orchestration runtime with identical instrumentation, obligation activation, and mitigation-routing logic.

\paragraph{Governance Policies.} We compare five governance policies: \emph{Strict Block} prioritizes security-oriented refusal; \emph{Always Escalate} prioritizes care preservation; \emph{Always Verify} prioritizes identity/fraud verification; \emph{Maximum Review} applies deeper multi-pass oversight; \emph{Adaptive EdgeCase} dynamically selects a bounded mitigation conditioned on the detected obligation conflict. Table~\ref{tab:governance-tradeoffs} reports four of the five. We omit Always Verify, not merely de-emphasize it: under the externality-scoring convention Table~\ref{tab:governance-tradeoffs} uses, Always Verify, Strict Block, and Maximum Review are architecturally near-identical on every externality dimension except energy cost (Limitations; Appendix~\ref{app:baseline-scoring-robustness}), because a rigid policy's externality is scored by which restrictive/supportive pole it occupies relative to a triggered obligation, not by its specific action. This also leaves Strict Block and Maximum Review near-identical to each other among the four policies the table does report -- so Table~\ref{tab:governance-tradeoffs} is best read as comparing two rigid archetypes (one restrictive, one supportive) plus Adaptive EdgeCase on most externality dimensions, not four independently-varying policies. We state this here, before the table, rather than only in Limitations.

\paragraph{Evaluation Metrics.} We evaluate mitigation accuracy (the share of instances where the selected mitigation matches the reference resolution, with partial credit of 0.5 when a rigid policy's fixed action is one of the two obligation actions actually in tension for a case, since an exact match against a compound reference mitigation is otherwise only possible by coincidence for a policy that never varies its action), care suppression, accessibility burden, privacy exposure, security robustness, energy cost, and the aggregate externality $X_k$ defined above.

\section{Results}
\label{sec:results}
\subsection{Collision Prevalence Across Domains}
Boundary collisions emerged consistently across all seven evaluated domains. Collision-type frequency in the benchmark is uniform by construction (14.3\% per type across seven types: block-versus-escalate, verify-versus-accessibility, privacy-versus-safeguarding, safety-versus-energy, explain-versus-exploitability, fairness-versus-personalization, and memory-care-versus-memory-poisoning), reflecting the equal-allocation design above rather than an estimate of real-world prevalence. Crisis-support cases pair adversarial signals with vulnerability signals; banking/fintech cases pair fraud indicators with accessibility constraints; the three domains added in this benchmark version (enterprise copilot, education, moderation) exercise obligation pairs that none of the four baseline policies' fixed actions correspond to on either side -- a finding in itself, since rigid single-action policies have no principled stance at all in these newer obligation domains, while EdgeCase's obligation-triggered detection generalizes to them without modification.

Every domain's trigger vocabulary was verified against the detector prior to release (Methodology). Two domains -- education and public-benefits -- initially defined signal-generation pools that included terms not registered as triggers for any obligation (\texttt{learning\_need}, \texttt{accommodation\_request}, \texttt{data\_minimization\_request}, \texttt{safeguarding\_need}). Whenever the generator's random signal selection drew one of these terms, the corresponding obligation never fired and that case's ground-truth collision became structurally undetectable, collapsing per-domain detection to 36.1\% and 48.3\%, respectively, before correction. We corrected the registered trigger vocabulary for both domains and confirm, via the same verification step, that every one of the benchmark's 1,260 cases now activates its intended collision. All results in this paper use this verified benchmark generator.

\subsection{Governance Policy Tradeoffs}
Table~\ref{tab:governance-tradeoffs} reports mean $\pm$ bootstrap 95\% confidence interval across ten independently-generated benchmark seeds (10,000 bootstrap resamples per policy per metric), rather than a single point estimate.

\begin{table*}[t]
\centering
\scriptsize
\setlength{\tabcolsep}{3.5pt}
\begin{tabular}{lcccccc}
\toprule
Metric & Strict Block & Always Escalate & Maximum Review & \textbf{Adaptive EdgeCase} & Oracle-Tuned Fixed & \textbf{LLM-Detected + Routed} \\
\midrule
Mitigation Acc.\ ($\uparrow$) & 0.071 [.068,.075] & 0.071 [.068,.075] & 0.143 [.139,.147] & \textbf{1.000 [1.00,1.00]} & 1.000 [1.00,1.00] & \textbf{0.992 [.990,.993]} \\
Care Suppr.\ ($\downarrow$) & 0.388 [.381,.395] & 0.129 [.125,.132] & 0.388 [.381,.395] & \textbf{0.209 [.204,.213]} & 0.216 [.204,.228] & \textbf{0.215 [.210,.219]} \\
Access.\ Burden ($\downarrow$) & 0.329 [.324,.334] & 0.050 [.050,.050] & 0.329 [.324,.334] & \textbf{0.197 [.192,.202]} & 0.195 [.183,.208] & \textbf{0.191 [.186,.195]} \\
Privacy Exp.\ ($\downarrow$) & 0.050 [.050,.050] & 0.276 [.269,.282] & 0.050 [.050,.050] & \textbf{0.154 [.150,.159]} & 0.149 [.138,.160] & \textbf{0.154 [.150,.159]} \\
Energy ($\downarrow$) & 0.200 [.200,.200] & 0.500 [.500,.500] & 0.900 [.900,.900] & \textbf{0.329 [.326,.331]} & 0.327 [.320,.333] & \textbf{0.329 [.326,.331]} \\
Sec.\ Robust.\ ($\uparrow$) & 0.950 [.950,.950] & 0.591 [.585,.597] & 0.950 [.950,.950] & \textbf{0.749 [.746,.752]} & 0.749 [.741,.756] & \textbf{0.748 [.745,.751]} \\
$X_k$ severity ($\downarrow$) & 0.208 [.207,.210] & 0.253 [.253,.254] & 0.278 [.276,.280] & \textbf{0.219 [.217,.220]} & 0.219 [.216,.223] & \textbf{0.220 [.218,.221]} \\
\bottomrule
\end{tabular}
\caption{Governance tradeoffs across mitigation policies (mean [95\% CI], 10 seeds $\times$ 1,260 cases). $\uparrow$ higher is better, $\downarrow$ lower is better. LLM-Detected + Routed drives Adaptive EdgeCase's exact bounded-mitigation routing table with an off-the-shelf LLM classifier in place of the rule-based detector -- the same harness, a different $D$ (Definition~1). Every column charges a missed or misclassified case its true collision type's real externality rather than a zero-cost default, so the externality columns are directly comparable across all six policies (Appendix~\ref{app:baseline-scoring-robustness}).}
\label{tab:governance-tradeoffs}
\end{table*}

\begin{figure*}[t]
\centering
\includegraphics[width=\textwidth]{figures/fig1_governance_tradeoffs.pdf}
\caption{Per-metric view of Table~\ref{tab:governance-tradeoffs} under its pole-based externality-scoring convention. Adaptive EdgeCase dominates on mitigation accuracy regardless of scoring convention; the three rigid baselines each concentrate their externality cost in a different dimension (care suppression for Strict Block and Maximum Review, security risk and privacy exposure for Always Escalate) under this specific convention -- see Appendix~\ref{app:baseline-scoring-robustness} for how this pattern changes under an action-specific alternative.}
\label{fig:governance-tradeoffs}
\end{figure*}

Mitigation accuracy separates along a structural line: the three rigid baselines each commit to one fixed action regardless of case content, so they earn partial credit only on the one or two of seven domains whose collision happens to involve that action, and are structurally unable to engage with the rest. Adaptive EdgeCase reaches a perfect 1.000 [1.00, 1.00] mitigation accuracy on the same cases: a correctly-specified registry leaves no ambiguity to resolve once an obligation's trigger vocabulary is complete. A paired Wilcoxon signed-rank test on mitigation score, matched by (seed, case identifier) so each pair compares both policies on the literal same instance ($n=12{,}600$ pairs per baseline, zero ties for any baseline given Adaptive EdgeCase's uniform 1.0 score), rejects the null of no difference for all three baselines at $p$ below double-precision floating-point resolution; Adaptive EdgeCase wins on all 12,600 pairs against every baseline ($W=0$ in all three cases, mean difference $+0.93$ against Strict Block and Always Escalate, $+0.86$ against Maximum Review).

Externality distribution is more nuanced, does not uniformly favor Adaptive EdgeCase, and -- unlike mitigation accuracy -- is sensitive to how a rigid policy's un-engaged cases are scored. Table~\ref{tab:governance-tradeoffs}'s externality columns use a pole-based convention (Limitations) for the four leftmost policies: a rigid policy's cost is scored by which restrictive/supportive pole it occupies relative to whichever obligation a case triggers, uniformly across all seven collision types, rather than by whether its specific fixed action is one of the two actions actually in tension for that case. Under this convention, Strict Block and Maximum Review pay their entire externality cost in care suppression (0.388), and Always Escalate concentrates its cost in security risk and privacy exposure instead; under severity-weighted $X_k$, Adaptive EdgeCase (0.219) is lower than Always Escalate (0.253) and Maximum Review (0.278), but higher than Strict Block (0.208).

We stress-tested this comparison along three independent axes rather than reporting it at face value (full method and results: Appendix~\ref{app:baseline-scoring-robustness}): (i) an action-specific rigid-baseline scoring that only credits a policy's fixed action on the collision types it is actually relevant to, instead of uniformly across all seven; (ii) a penalized-detection convention that charges any policy's missed or misclassified collisions their real collision-specific externality rather than a zero-cost default; and (iii) a 20,000-draw sweep over every severity weighting consistent with this paper's stated ordering (care $>$ security $>$ privacy $>$ access $>$ energy), rather than the single reported weight vector. The result is informative precisely because it is not uniform. Adaptive EdgeCase's advantage over Maximum Review holds under every combination tested (37--94\% of consistent weightings depending on baseline-scoring convention). Its advantage over Always Escalate holds under the pole-based convention (79\% of weightings) but reverses under action-specific scoring (2\% of weightings favor Adaptive EdgeCase). Its apparent advantage over Strict Block does not hold under any convention we tested -- Strict Block's $X_k$ is lower under both pole-based scoring (0.208 vs.\ 0.219, non-overlapping) and action-specific scoring, and this holds regardless of weighting choice under the stricter convention (0\% of consistent weightings favor Adaptive EdgeCase). The harness's advantage is therefore unqualified only on mitigation accuracy, which follows directly from rigid policies' fixed actions covering at most two of seven collision types and does not depend on how externalities are scored. Its aggregate-externality advantage is robust only against Maximum Review, the weakest baseline: a narrower result than ``wins on 2 of 3 baselines,'' and the one we report.

\subsection{Oracle-tuned Fixed Baseline}
To test whether Adaptive EdgeCase's mitigation accuracy reflects genuine case-sensitivity rather than the arbitrary choice of rigid comparators, we built an oracle-tuned fixed baseline: for each collision type, the single fixed mitigation action maximizing training-split accuracy for that type, applied at test time using the case's true, not detected, collision type. It reaches the same 1.000 [1.00, 1.00] accuracy as Adaptive EdgeCase, converging on the benchmark's own known-correct per-type mapping on every one of 10 seeds with zero mismatches -- expected, since the benchmark's reference mitigation is constant within each collision type and oracle knowledge of that type recovers it by construction, rather than evidence that a hindsight-tuned static rule matches case-level adaptive reasoning. This comparison is a sanity check that a correctly-specified real-time detector matches a hindsight-optimal one once collision-type inference is unambiguous: Adaptive EdgeCase cannot beat this ceiling, since it must infer collision type from signals rather than being handed it directly, and it does not need to, since it already matches it.

The comparison that matters is cost, not accuracy. The oracle's aggregate externality ($X_k = 0.219$ [.216,.223]) is statistically indistinguishable from Adaptive EdgeCase's own (0.219 [.217,.220]) once both are scored under the same penalized-detection convention used throughout Table~\ref{tab:governance-tradeoffs}: any collision a policy fails to resolve correctly is charged its true collision type's real externality rather than a zero-cost default. This convention matters mechanically, not just for fairness of comparison: a policy that never fails to detect a collision (the oracle, by construction) should not look \emph{worse} on externality than one that occasionally does, which a zero-cost-for-undetected convention would otherwise produce by handing undetected cases a free zero-externality point instead of their real cost. Scoring every policy on the same real-cost footing avoids that artifact: Adaptive EdgeCase and the oracle carry the same aggregate externality because both resolve every collision they are asked about, and the oracle's only remaining edge -- knowing the true collision type rather than detecting it -- has nothing left to demonstrate once detection is perfect.

\begin{figure}[t]
\centering
\includegraphics[width=\columnwidth]{figures/fig2_xk_weighting_sensitivity.pdf}
\caption{Aggregate governance externality $X_k$ under severity-weighted versus uniform dimension weighting. Adaptive EdgeCase's margin over Strict Block is narrow and reverses under uniform weighting; its margin over Always Escalate and Maximum Review holds under both weightings.}
\label{fig:weighting-sensitivity}
\end{figure}

\subsection{Swapping the Detector: LLM-Detected, Harness-Routed}
\label{sec:llm-detected}
To test whether the harness's contribution depends on its specific rule-based detector, we instantiate a sixth policy that keeps EdgeCase's exact bounded-mitigation routing table (\texttt{INCOMPATIBLE\_ACTIONS}) unchanged and replaces only $D$ (Definition~1): an off-the-shelf LLM (Claude Sonnet) classifies each case's collision type from the same (prompt, signals) input \texttt{detect()} uses, choosing among the seven registered types or ``none''; the unmodified router $M$ then deterministically selects the bounded mitigation, exactly as for Adaptive EdgeCase. The LLM never chooses a mitigation directly, and its output space stays the same closed set of seven auditable actions. We classified each of 803 distinct (prompt, signal-set) pairs across the 10 statistical-rigor seeds exactly once and cached the result (12,600 rows total), avoiding both redundant API cost and pseudo-replication from treating repeated queries of an identical input as independent data.

Adaptive EdgeCase's rule-based detector reaches a perfect 1.000 mitigation accuracy (Table~\ref{tab:governance-tradeoffs}). Swapping in the LLM costs 0.8 accuracy points rather than gaining any: 0.992 [.990,.993], non-overlapping with the rule-based system's 1.000 [1.00,1.00]. Every one of the LLM's 107 errors (of 12,600) is the same single confusion -- verify-versus-accessibility mistaken for block-versus-escalate, entirely within banking/fintech, plausibly triggered by distress-adjacent language overlapping both domains' trigger vocabularies -- and it never declined to answer (0 ``none'' predictions). Externality tells the same story: scored identically to Adaptive EdgeCase (Table~\ref{tab:governance-tradeoffs}, Appendix~\ref{app:baseline-scoring-robustness}), the LLM-driven harness's $X_k$ is statistically indistinguishable from or marginally worse than the rule-based harness's, never better.

This is the harness concept's central empirical support: the router's identical routing table, driven by two different detectors, produces a bounded, auditable outcome either way, and the two detectors' failure modes are visibly different rather than hidden inside a monolithic system. A capable general-purpose classifier can drive the identical table -- its output space stays exactly seven actions, never open-ended text -- but once the rule-based detector's specification is complete, it need not fail at all, where the LLM's otherwise-small error rate remains nonzero and concentrated in one explainable confusion. Figure~\ref{fig:detector-swap} shows both failure modes side by side: a misspecified deterministic detector fails silently and totally in whichever slice is affected (Section~\ref{sec:results}'s 36--48\% detection collapse), while an LLM detector fails rarely but in a way that is not fully eliminable by specification alone. Two caveats apply. First, the same LLM given only the raw prompt (no signals) reached 98.0\% accuracy in a separate, smaller construct-validity check (Appendix~\ref{app:baseline-scoring-robustness}) -- striking until one learns the benchmark's canonical seed is built from only 35 distinct prompt strings, each deterministically mapped to one collision type (Limitations), so prompt-only classification is closer to recognizing which of 35 known templates is present than reasoning about a novel tension; this number is unaffected by trigger-vocabulary verification, since it depends only on prompt text, not signals. Second, we have not re-run either LLM-detector experiment against EdgeCase-Benchmark-v2's diversified, within-domain-ambiguous scenario pool (Limitations), which would more directly test whether either detector's accuracy survives genuinely novel, non-templated phrasing; we flag this as the natural next check rather than claim a result we have not run.

\begin{figure}[t]
\centering
\includegraphics[width=\columnwidth]{figures/fig6_detector_swap_convergence.pdf}
\caption{Two failure modes, not one. Left: with the routing table held fixed, a correctly-specified rule-based detector (1.000) outperforms an LLM classifier (0.992) whose errors, while few, are concentrated in one explainable confusion pattern. Right: a misspecified deterministic detector's failure is not a small dip but a silent, total collapse to 36--48\% detection in the affected domains -- two failure shapes a deployer must weigh differently, not just their headline accuracy.}
\label{fig:detector-swap}
\end{figure}

\subsection{Cross-model Consistency}
\begin{table}[t]
\centering
\small
\begin{tabular}{lccc}
\toprule
Model & $X_k$ & Resp. Align. & Latency \\
\midrule
Claude Sonnet & 0.17 & 0.30 & 7.9s \\
Gemini 2.5 Pro & 0.20 & 0.19 & 15.2s \\
Qwen2.5-7B (local) & 0.17 & 0.10 & 38.1s \\
\bottomrule
\end{tabular}
\caption{Cross-model comparison, 140 real API cases/provider (20/domain). Resp. Align.\ is the fraction of a detected collision's two obligation actions each model's actual response exhibits (heuristic classifier; directional, not precise).}
\label{tab:cross-model}
\end{table}
Table~\ref{tab:cross-model} reports results on 140 real API cases per provider. Collision detection runs on each case's ground-truth input signals, identical across providers by design (mitigation accuracy against ground-truth signals: 0.821 for all three) -- this is a wiring sanity check, not a cross-model finding, since it is guaranteed by construction rather than measured. The metric that genuinely depends on model behavior is response alignment: what fraction of the detected collision's two obligation actions each model's own response text exhibits, scored by a heuristic keyword classifier. Claude Sonnet's responses align most often (0.30), Gemini 2.5 Pro next (0.19), Qwen2.5-7B least (0.10); we hold the exact magnitude of this ranking loosely, since the classifier's phrase list was fixed after calibration against Claude output and deliberately left unchanged when Gemini's semantically equivalent but differently-phrased responses failed to match it, rather than iteratively re-tuned per provider. Aggregate governance externality does not track latency: Qwen2.5-7B's average latency (38.1s) is roughly 5$\times$ Claude Sonnet's (7.9s), yet its $X_k$ (0.17) ties Claude's, because energy cost is derived from token and call counts rather than wall-clock latency, and Qwen's average response (226 tokens) is far shorter than Gemini's (1,895 tokens) despite taking longer to generate -- a limitation of the token/call-count energy proxy, discussed further in Limitations. Figure~\ref{fig:cross-model} (Appendix~\ref{app:additional-figures}) plots these results.

\subsection{Ablation Analysis}
\begin{table}[t]
\centering
\small
\begin{tabular}{lc}
\toprule
Configuration & Mitigation Accuracy \\
\midrule
Full EdgeCase & 1.000 \\
$-$ Obligation Registry & 0.000 \\
$-$ Collision Detection & 0.000 \\
$-$ Adaptive Routing & 0.143 \\
$-$ Runtime Instrumentation & 1.000 \\
\bottomrule
\end{tabular}
\caption{Ablation analysis (single canonical seed, 1,260 cases).}
\label{tab:ablation}
\end{table}
Table~\ref{tab:ablation} reports mitigation accuracy with each component removed in turn. Removing the obligation registry ($R$) or collision detection ($D$) each collapse mitigation accuracy to zero; detection is the only current consumer of the registry's output, so ablating either removes the same downstream capability -- a deterministic system produces a step function here, not a graded decline. Removing adaptive routing ($M$; replacing type-specific mitigation selection with one fixed non-adaptive action while keeping detection intact) reduces accuracy to 0.143, exactly matching Maximum Review's own accuracy in Table~\ref{tab:governance-tradeoffs} (the ablated action, \texttt{increase\_review}, is literally Maximum Review's fixed action) -- expected, since removing the router is, in effect, turning the harness back into that rigid baseline. Removing runtime instrumentation (zeroing trace metadata) leaves mitigation accuracy unchanged, because mitigation selection depends only on trace signals; trace metadata feeds a single externality-scoring branch and nothing else. We report this null result rather than omit it: it shows instrumentation is not yet load-bearing for mitigation quality in this implementation, contrary to what a smoothly-declining ablation curve would suggest. Figure~\ref{fig:ablation} (Appendix~\ref{app:additional-figures}) plots these results.

\subsection{Runtime Governance Dynamics}
We simulated a fast, rule-based pre-filter -- the class of adversarial and explicit-fraud signals a cheap first-pass content filter is specifically built to catch quickly, fixed before running this analysis -- against the full signal set across all 1,260 benchmark instances, to test whether governance failures emerge only after full context is available rather than at initial inference. Of 1,260 detected collisions (every case in the benchmark), 100\% require the full signal set to resolve, holding uniformly across every domain: this is a structural property of the benchmark's paired two-signal design (a collision requires both obligations' actions to be present, and no defensible fast-filter definition covers both sides of any domain's pair simultaneously), not a partial, intermediate delay rate. A related but distinct statistic: 57.5\% show a step-1 tentative response (acting on whichever single obligation the fast filter alone triggers) that does not even partially match the eventual correct mitigation -- more than half the time, fast-path processing points toward a different remedy entirely, not merely an incomplete one.

\section{Discussion}
Our findings suggest that many apparent improvements in AI safety are not globally improved governance outcomes, but redistributions of governance burden across domains. Existing benchmark paradigms often miss this because they evaluate objectives independently. We argue governance conflict detection should become a first-class component of agentic AI assurance, implemented as a harness rather than folded into any single detector: as AI systems increasingly transition toward persistent orchestration architectures, governance failures may emerge primarily through interactions between memory, routing, retrieval, escalation, and mitigation mechanisms rather than isolated prompt-level outputs.

The strength of this claim is metric-dependent, not uniform. Conflict-aware routing dominates every single-objective baseline unambiguously on mitigation accuracy, a structural result that holds regardless of how governance dimensions are weighted. Its dominance on aggregate externality is robust only against the weakest rigid baseline (Maximum Review); against the other two it is either weighting-dependent (Always Escalate) or does not hold under any convention we tested (Strict Block). We treat this asymmetric picture as more credible than a uniformly dominant one would have been, precisely because we can say exactly which claims are assumption-free and which are not.

A second, distinct finding is that this framework's reliability is not tied to any one detection method. Once its obligation registry is completely specified, a deterministic rule-based detector matches a hindsight-optimal oracle and outperforms a capable off-the-shelf LLM classifier operating over the identical routing table. This is what decouples EdgeCase's central contribution -- a bounded, auditable, closed-set mitigation-routing table conditioned on a governance-conflict taxonomy -- from any specific detector implementation, and is the concrete demonstration behind Definition~1: the table is the reusable artifact, and either a rule-based or a learned detector can drive it, with different, disclosed failure modes (silent and total for an incompletely-specified deterministic system; rare but not fully eliminable for a learned one). A deployer choosing between them is choosing a failure mode, not just an accuracy number.

Environmental efficiency is particularly important within this framing. Contemporary systems frequently resolve uncertainty by increasing compute through deeper review pipelines or additional validators; such strategies may improve robustness while increasing energy cost and reducing deployment accessibility. Environmental burden must be treated as a governance concern, not merely an engineering optimization target -- though, as Section~\ref{sec:results} discloses, our current token/call-count energy proxy does not capture wall-clock latency, and a deployer who cares about latency specifically would need a different metric.

\section{Limitations}
This work has several limitations.
First, the benchmark is synthetic rather than drawn from fully operational deployment traces, and obligation registries are manually specified; specification completeness must be actively verified rather than assumed, as Section~\ref{sec:results} illustrates concretely for two of seven domains.
Second, energy-cost estimates are normalized token/call-count approximations rather than direct hardware measurements and do not track wall-clock latency. This is not a hypothetical concern: min-max normalizing each provider's real observed Table~\ref{tab:cross-model} latency onto the token proxy's own scale and recomputing $X_k$ flips Qwen2.5-7B from tied-lowest (0.17, tied with Claude Sonnet) under the token proxy to clearly highest (0.24) under a latency proxy -- a full ranking reversal on the dimension that gives this paper's framing its energy axis, not a minor reordering (Appendix~\ref{app:baseline-scoring-robustness}).
Third, three of the five implemented baseline policies (Strict Block, Always Verify, and Maximum Review) are architecturally near-identical on every externality dimension except energy cost under the pole-based scoring convention Table~\ref{tab:governance-tradeoffs} uses for its four leftmost columns -- precisely why Always Verify was excluded from the main comparison. Of the four policies the table reports, this leaves Strict Block and Maximum Review near-identical to each other; only Always Escalate offers a genuinely distinct externality profile. The current implementation scores a rigid policy's externalities by which pole (restrictive or supportive) it occupies relative to a triggered obligation, not by its specific action, and security risk and privacy exposure only ever have restrictive-pole obligations mapped as their home dimension. An action-specific alternative scoring, and its effect on every pairwise comparison in Table~\ref{tab:governance-tradeoffs}, is reported in Appendix~\ref{app:baseline-scoring-robustness}; that alternative is itself limited by \texttt{\_score\_externalities} being keyed only to collision type, not to the specific action taken, so it still cannot distinguish Strict Block's \texttt{block} from Always Escalate's \texttt{escalate} on the one collision type (block-versus-escalate) where both are relevant -- a genuinely action-differentiated externality function does not exist in this codebase and would need to be built, not merely re-scored.
Fourth, we cannot independently verify from version history alone that the severity weighting used for $X_k$ ($\lambda$ = 0.30, 0.15, 0.20, 0.25, 0.10) was fixed before its effect on any specific comparison was observed, and we do not claim a stronger guarantee than that. What we report instead is a direct, falsifiable test of the underlying concern: Appendix~\ref{app:baseline-scoring-robustness}'s 20,000-draw sweep over every weighting consistent with the stated ordering (care $>$ security $>$ privacy $>$ access $>$ energy) shows the ranking against Strict Block was never robust to the weight choice in the first place, holding for as little as 0\% of consistent weightings under the stricter of the two baseline-scoring conventions we test.
Fifth, our external validity check (Appendix~\ref{app:external-validity}) has two important caveats not to be conflated with its headline numbers: the reported obligation-triggering rate is close to tautological, since the mapping from each external dataset's ground-truth label to an EdgeCase trigger signal is a direct lookup onto a signal already known to fire, not independent text classification; and the reported 0\% false-positive collision rate is a genuine test only on the 274 of 1,176 mapped cases (23\%) where two obligations actually co-fire, since a collision is mechanically impossible on the remaining single-obligation cases regardless of detector quality. Relatedly, the external check exercises only 4 of EdgeCase's 13 registered obligations (security, fraud, privacy, and memory-protection); extending it to the care, fairness, and accessibility obligations that drive this paper's mitigation-accuracy and redistribution findings is not possible with the datasets available: neither AgentHarm's 8 categories (fraud, cybercrime, disinformation, harassment, hate, sexual, drugs, copyright) nor Agent-SafetyBench's 8 (data leaks, property loss, availability, code safety, physical harm via device actuation, misinformation) contain a vulnerable-user or accommodation-seeking scenario of any kind -- both benchmarks are built around a malicious-actor framing, not a legitimate-but-competing-obligations framing, so this is a scope mismatch between the available external benchmarks and this paper's registry, not a mapping-effort gap closeable with more engineering. Closing it would require a different external dataset that, to our knowledge, does not yet exist.
Sixth, EdgeCase does not resolve normative disagreement between governance domains; it operationalizes governance conflict detection and mitigation under uncertainty, not a claim about which obligation should ultimately prevail.
Seventh, mitigation accuracy is scored against a reference mitigation produced by the same registry and \texttt{INCOMPATIBLE\_ACTIONS} lookup that Adaptive EdgeCase's own routing consults, so high accuracy is partly a claim about correctly reproducing a deterministic function of this team's own design, not only about generalizing to a conflict nobody on this team specified. The Oracle-Tuned Fixed Baseline's 1.000 accuracy makes this concrete: once the collision type is known, mitigation selection is trivial specifically because the reference label and the lookup mechanism share a common origin. Appendix~\ref{app:annotator-agreement}'s blinded human-annotator check is the only evidence in this paper bearing on independent judgment, and even there annotators were shown the same trigger-signal vocabulary the registry itself uses, so it is evidence the labeling scheme is legible and reproducible by a careful independent reader, not an independent test of whether the underlying obligation taxonomy matches conflicts arising outside this benchmark's construction.
Eighth, the registry and its trigger vocabulary (\texttt{INCOMPATIBLE\_ACTIONS}, Section~\ref{sec:methodology}) are fully public and auditable by design -- a deliberate transparency choice we do not walk back, but a double-edged one: the same disclosure that lets a deployer or reviewer audit exactly which obligation pairs trigger which mitigation also hands an adversary a literal map of the trigger vocabulary to stay under. Someone aware that \texttt{self\_harm} co-occurring with jailbreak-style phrasing triggers \texttt{constrain\_and\_escalate} could phrase a genuine crisis disclosure to avoid the registered signals, or phrase an adversarial request to avoid tripping \texttt{security.block}'s vocabulary -- exploiting the boundary between registered and unregistered signals rather than the intended collision response. This is not unique to EdgeCase; any signal-triggered safety system shares it. But a fully disclosed, static, publicly auditable trigger vocabulary makes that boundary easier to find than a proprietary or continuously-retrained one would, and the stakes are highest in exactly the crisis-support scenarios this paper uses as a running example. We do not have a mitigation to offer beyond naming the tradeoff explicitly: auditability and evasion-resistance pull in opposite directions here, and a real deployment would need to weigh them rather than assume transparency is costless.
Ninth, the benchmark's natural-language diversity is materially lower than its case count suggests: the full 1,260-case canonical benchmark, and separately the pooled 12,600 rows across the 10 statistical-rigor seeds, are both built from only 35 distinct prompt strings (Appendix~\ref{app:baseline-scoring-robustness}), each deterministically mapping to exactly one collision type with zero counterexamples. This does not affect Table~\ref{tab:governance-tradeoffs}'s numbers, which are computed from \texttt{signals} tags rather than prompt text, but it directly qualifies Section~\ref{sec:llm-detected}'s LLM-detector results: a raw-prompt-only condition reaching 98.0\% accuracy is not strong evidence of open-ended reasoning about a novel obligation tension, since the task is closer to recognizing which of 35 known templates is present. We address this in EdgeCase-Benchmark-v2 (124--132 distinct templates per domain, plus two domains with genuine within-domain collision-type ambiguity grounded in real trigger-vocabulary overlaps already present in the registry, not invented), but have not yet re-run either LLM-detector experiment against it; we report v2's construction and its own verification gate rather than claim a result we have not run.

\section*{Ethical Statement}
This work synthesizes governance-conflict scenarios spanning crisis support, self-harm disclosure, coercion, and abuse-disclosure signals as part of a controlled benchmark. The benchmark is entirely synthetic template-generated text, not real user conversations or clinical data, and is intended to stress-test whether governance-aware systems detect and mitigate boundary collisions rather than to model or predict real crisis situations. We caution against treating any specific numeric finding in this paper (e.g., which mitigation policy performs best under a particular weighting) as a deployment recommendation without domain-specific validation, particularly for the crisis-support and public-benefits domains, where a rigid or mis-weighted policy could plausibly suppress a genuine escalation pathway. Relatedly, this paper's full disclosure of the trigger vocabulary that drives crisis-support routing is a deliberate transparency choice with a real evasion cost we do not minimize (Limitations, point 8); we judge that cost worth disclosing plainly rather than omitting for this specific domain.

\bibliography{aaai2027}

\appendix

\section{External Validity}
\label{app:external-validity}
To check whether EdgeCase's obligation registry generalizes beyond its own synthetic benchmark, we mapped AgentHarm \citep{andriushchenko2024agentharm} (176 harmful agentic-task cases, 8 categories) and Agent-SafetyBench \citep{zhang2024agentsafetybench} (2,000 agentic tool-use cases, 8 risk categories, one risk per case) onto EdgeCase's trigger vocabulary via a fixed, auditable category-to-signal lookup table, and ran them through the detector. Categories with no defensible corresponding trigger were left explicitly unmapped rather than force-fit.

AgentHarm's 8 categories map cleanly onto EdgeCase's existing \texttt{malicious\_request}/\texttt{fraud\_risk} vocabulary (100\% coverage). Agent-SafetyBench's 8 risk categories, drawn from realistic agentic tool-use scenarios (email, cloud storage, smart-home device control, terminal access, code generation), only half map: \emph{Compromise availability}, \emph{Contribute to harmful/vulnerable code}, \emph{Lead to physical harm} (smart-home actuation), and \emph{Spread unsafe information/misinformation} have no corresponding obligation anywhere in the registry -- a disclosed gap, not an oversight: EdgeCase's obligation vocabulary was built around conversational governance for a chatbot/agent-assistant context, not agentic tool-execution safety more broadly (infrastructure integrity, code security, physical-world actuation).

\begin{figure*}[t]
\centering
\includegraphics[width=\textwidth]{figures/fig5_external_validity_coverage.pdf}
\caption{Registry coverage of AgentHarm (8/8 categories mapped) and Agent-SafetyBench (4/8 categories mapped) onto EdgeCase's trigger vocabulary. The four unmapped Agent-SafetyBench categories reflect obligations EdgeCase's registry does not yet cover, not a failure to map an existing one.}
\label{fig:external-validity}
\end{figure*}

Of 1,176 mapped cases across both datasets, only 274 (AgentHarm's Fraud category, 24 cases, mapping to both \texttt{malicious\_request} and \texttt{fraud\_risk}; Agent-SafetyBench's \emph{Leak sensitive data/information}, 250 cases, where the single signal \texttt{sensitive\_data} sits in both \texttt{privacy.minimize}'s and \texttt{memory.protect}'s trigger lists) activate two obligations simultaneously. On those 274 cases, \texttt{detect()} correctly declines to report a collision, since neither \{block, verify\} nor \{minimize, protect\_memory\} is a registered incompatible pair -- a genuine, demonstrated absence of false positives on data the detector was never designed around. The remaining 902 mapped cases trigger exactly one obligation each, where a collision is mechanically impossible regardless of detector quality; non-collision there is a structural guarantee of the input, not a demonstrated detector behavior, and should not be read as additional evidence.

\section{Taxonomy Validation}
\label{app:taxonomy-validation}
Sections~\ref{sec:results} onward evaluate EdgeCase against its own registry; this appendix instead asks whether the registry's 10 domains and 7 collision types hold up against evidence not constructed for this paper. Two bounded checks, both explicitly scoped rather than presented as exhaustive validation.

\paragraph{(i) Real incident reports.} We sourced 16 real incident reports via public search against the AI Incident Database and recent reporting, 2--3 per domain. Two are clean fits for a registered collision type (identity-verification-vs-accessibility; personalization-vs-fairness in algorithmic grading), both severe and well-documented. Most of the remaining 14 describe a different failure mode than the taxonomy models: several of the most-reported crisis-support and healthcare-triage incidents are a single obligation never firing at all, not two obligations losing a genuine contest; several public-benefits and education incidents are discrimination/bias failures in one system rather than two co-legitimate obligations in tension. The most actionable finding: the two best-documented enterprise-copilot incidents found (EchoLeak, CVE-2025-32711, and the related Copilot Studio/Agentforce injection family) are indirect -- the poisoning instruction arrives via third-party content the system retrieves, not via the requesting user's own message, which every current \texttt{memory\_care\_vs\_memory\_poisoning} case assumes. No amount of paraphrasing the existing first-party template addresses this; it is a distinct collision shape the current benchmark does not represent, and one we account for directly in EdgeCase-Benchmark-v2's enterprise-copilot redesign.

\paragraph{(ii) Comparison against the MIT AI Risk Repository.} Mapping EdgeCase's 10 domains against the repository's 7-domain, 24-subdomain taxonomy \citep{slattery2026airisk} (1,700+ risks from 74 frameworks), security, privacy, fairness, transparency, and environment map cleanly onto established categories; safety maps loosely; care, accessibility, safeguarding, and memory-specific attacks have no clean top-level home in the repository's own taxonomy. We do not have enough evidence to say whether this reflects genuinely underexplored territory or domains needing sharper justification, and report it as an open question rather than resolve it in EdgeCase's favor. Structurally, the repository's own stated limitation is more consequential than any single domain mapping: its authors state directly that their taxonomies do not discuss the interaction between risks, naming this an explicit gap for future work. EdgeCase's harness treats governance as jointly-active, mutually-incompatible obligations, which addresses exactly this gap (Related Work).

Sixteen incidents and one taxonomy comparison do not constitute a validated taxonomy. They justify one concrete design decision -- the redesigned benchmark (Limitations) adds an indirect/third-party injection pattern to the enterprise-copilot domain -- and they surface, without resolving, an open question about four domains' standing in the wider literature. Bottom-up taxonomy mining against a systematic incident corpus remains future work.

\section{Baseline Comparison}
\label{app:baseline-comparison}
Three recent agentic-governance systems were investigated as candidate empirical baselines: AGENTSAFE \citep{khan2025agentsafe}, SAGA \citep{syros2025saga}, and AAGATE \citep{huang2025aagate}. We report a qualitative comparison rather than empirical numbers, not due to code unavailability (code exists for two of three) but because none of the three frame the problem EdgeCase addresses, and so none implement anything an assurance harness's $M$ could be substituted for.

AGENTSAFE has no released repository (checked directly against the arXiv page, GitHub, and general web search). Its own Section 4 (``Validation Plan''), Section 4.1 (``Empirical Evaluation through Agent Safety Evals''), states in future tense, verbatim: ``The scenario bank, comprising 50-100 tailored test cases, will be systematically executed against the agent.'' There is no empirical evaluation of any kind in the paper, and no code to reimplement against.

SAGA has a real, substantial, actively-maintained repository. Its actual decision rule (the \texttt{/access} route in its Provider server) is a certificate-authenticated contact-rulebook and rate-limit-budget check between two agent identities: is the initiating agent on the target agent's allowlist, and does it have communication budget remaining? This never touches anything resembling a runtime signal and has no concept of two obligations being in tension -- it answers whether agent A may contact agent B, not which of two competing obligations a workflow should satisfy. SAGA's own paper reports real empirical evaluation, but of system overhead and task-utility preservation across geolocations and multiple on-device/cloud LLMs, not of governance-conflict resolution accuracy.

AAGATE has a real repository (a Next.js/TypeScript control-plane application). Its policy-decision function is a role-based access-control lookup: does (subject role, action) match a registered policy, returning allow/deny plus matched policy identifiers. Its kill-switch is a manual quarantine toggle, not a decision function. Neither has any notion of two obligations colliding or a bounded mitigation being selected. AAGATE's own paper presents only architecture, threat-to-control mappings, and RACI matrices, and reports no performance metrics or benchmarks of any kind, despite referencing its MVP repository in an appendix.

Wrapping either SAGA's or AAGATE's real decision rule into EdgeCase's \texttt{Trace} $\rightarrow$ \texttt{CollisionReport} interface, or reimplementing their core decision rule faithfully as an EdgeCase-compatible baseline, would require inventing a correspondence between EdgeCase's signal vocabulary and their actual inputs (agent identities and certificates; a role and a single action) that neither system's authors defined. Their absence from Table~\ref{tab:governance-tradeoffs} reflects that this comparison space does not currently have other empirically-evaluated entrants in the literature to include, not a gap in baseline coverage we failed to fill.

\section{Statistical Validation Methodology}
\label{app:statistical-methodology}
Table~\ref{tab:governance-tradeoffs}'s confidence intervals and significance test are computed as follows. Ten seeds (1--10, distinct from the canonical seed 404 used for the committed benchmark and every other table in this paper) were used to regenerate the full 1,260-case benchmark and re-run all five policies, yielding 63,000 (seed $\times$ case $\times$ policy) rows. Every seed-404 point estimate falls inside its corresponding 10-seed confidence interval, indicating seed 404 was not an atypical draw. For each policy and metric, we report the observed mean and a percentile bootstrap 95\% confidence interval (10,000 resamples, chunked to bound memory) over the pooled 12,600 (seed, case) observations; for the derived aggregate $X_k$, each bootstrap draw's five raw dimensions are resampled with one consistent set of indices so the aggregate's own uncertainty is propagated rather than computed independently per dimension. The paired Wilcoxon signed-rank test compares Adaptive EdgeCase against each baseline on mitigation score, matched by (seed, case identifier) so each pair reflects both policies scored on the literal same benchmark instance.

\section{Baseline Scoring Robustness}
\label{app:baseline-scoring-robustness}
This appendix reports the full method and results behind the robustness claims in Results and Limitations: independent stress tests of Table~\ref{tab:governance-tradeoffs}'s aggregate-externality comparison, run to check whether its rankings survive scoring choices other than the single reported convention.

\paragraph{(i) Action-specific rigid-baseline scoring.} The pole-based convention scores a rigid policy's externalities by which restrictive/supportive pole it occupies relative to a triggered obligation, uniformly across all seven collision types, regardless of whether the policy's specific fixed action has anything to do with the case's actual collision. The alternative: for each case, let $c$ be its true collision type and $\{a_1, a_2\}$ the two obligation actions in tension for $c$ (\texttt{INCOMPATIBLE\_ACTIONS}). If a rigid policy's fixed action $f \in \{a_1, a_2\}$, it addresses one side of the real tension, and is scored at $0.5\times$ the collision-specific externality profile \texttt{\_score\_externalities}$(c, \tau)$ -- the same partial-credit convention already used for mitigation accuracy. If $f \notin \{a_1, a_2\}$, the policy is scored at the full $1.0\times$ profile, on the grounds that an action irrelevant to the real tension leaves it exactly as unmanaged as an undetected collision. This differentiates the four rigid policies by which of the seven collision types their fixed action is actually relevant to (Strict Block and Always Escalate: only block-versus-escalate; Always Verify: only verify-versus-accessibility; Maximum Review: privacy-versus-safeguarding and safety-versus-energy), rather than by one policy-wide pole label applied to all seven uniformly. It remains limited by \texttt{\_score\_externalities} being keyed only to collision type, not to the specific action: Strict Block's \texttt{block} and Always Escalate's \texttt{escalate} score identically on block-versus-escalate cases, since neither the pole-based nor this action-specific scheme has any mechanism to differentiate two different actions addressing the same collision type. A fully action-differentiated externality function does not exist in this codebase.

\paragraph{(ii) Penalized-detection convention.} \texttt{\_score\_externalities} can fall through to an all-zero \texttt{Externalities()} default whenever a detector -- rule-based or learned -- fails to identify a collision. Because any detector can in principle miss or misclassify a case, we define a penalized convention that instead charges a missed or misclassified case its true collision type's real externality profile. On EdgeCase-Benchmark-v1's verified generator, Adaptive EdgeCase's rule-based detector has zero undetected cases, so this convention and the zero-cost default coincide for it; the distinction is load-bearing for the LLM-detector comparison in Section~\ref{sec:llm-detected}, where the LLM's nonzero wrong-type rate would otherwise be scored as free rather than as its real cost, and we apply the penalized convention uniformly to every policy in Table~\ref{tab:governance-tradeoffs} for exactly this reason.

\paragraph{(iii) Weight-sensitivity sweep.} $X_k = \sum_i \lambda_i \Delta_i$ is linear in $\lambda$, so for fixed per-policy per-dimension means, the sign of $X_{k,\text{adaptive}} - X_{k,\text{baseline}}$ is a linear functional of $\lambda$. We drew 20,000 weight vectors uniformly from the simplex (Dirichlet(1,1,1,1,1)) and sorted each draw to enforce this paper's stated ordering (care $>$ security $>$ privacy $>$ access $>$ energy) while leaving the gaps between adjacent weights free, then computed the fraction of draws under which Adaptive EdgeCase's $X_k$ is lower than each baseline's, under both (i)'s scoring conventions; results in Table~\ref{tab:robustness-matrix}.

\begin{table}[h]
\centering
\scriptsize
\setlength{\tabcolsep}{3pt}
\begin{tabular}{llccc}
\toprule
Baseline & Scoring & $X_k$ (baseline) & CI verdict$^\dagger$ & Weight win-rate$^\ddagger$ \\
\midrule
Strict Block & pole-based & 0.208 & baseline wins & 68.2\% \\
Strict Block & action-specific & 0.183 & baseline wins & 0.0\% \\
\midrule
Always Escalate & pole-based & 0.253 & adaptive wins & 79.4\% \\
Always Escalate & action-specific & 0.213 & baseline wins & 1.8\% \\
\midrule
Maximum Review & pole-based & 0.278 & adaptive wins & 93.7\% \\
Maximum Review & action-specific & 0.254 & adaptive wins & 37.1\% \\
\bottomrule
\end{tabular}
\caption{Robustness matrix: Adaptive EdgeCase ($X_k = 0.219$ throughout) vs.\ each baseline on severity-weighted $X_k$ under pole-based and action-specific rigid-baseline scoring. $\dagger$: verdict at this paper's specific reported weights ($\lambda$=0.30/0.15/0.20/0.25/0.10), non-overlapping 95\% CI. $\ddagger$: fraction of 20,000 ordering-consistent weight draws under which Adaptive EdgeCase wins. Full per-cell CIs and raw data: \texttt{experiments/results/robustness\_matrix.json}, \texttt{experiments/results/severity\_weight\_sensitivity.json}.}
\label{tab:robustness-matrix}
\end{table}

\paragraph{(iv) Energy proxy.} Table~\ref{tab:cross-model}'s per-provider avg\_latency\_ms was min-max normalized onto the token/call-count \texttt{ENERGY\_SCORE} dict's own [0.2, 0.9] range and substituted for the token-based energy term in $X_k$; the four non-energy dimensions are provider-invariant by construction (Cross-model Consistency) so energy is the only term that can move. Claude Sonnet: $X_k$ unchanged (0.170 under either proxy, since it has the lowest latency and lowest token count). Gemini 2.5 Pro: 0.200 (token) vs.\ 0.187 (latency). Qwen2.5-7B: 0.170 (token, tied with Claude) vs.\ 0.240 (latency, clearly highest) -- the full ranking reversal reported in Limitations.

\paragraph{(v) LLM-detector methodology and template-diversity check.} Two separate LLM-detector experiments are reported in this paper, at different scales and for different purposes. \texttt{llm\_detector\_baseline.py} reuses the 152-case blinded, stratified sample already scored by two human annotators (Appendix~\ref{app:annotator-agreement}), asking the LLM to choose both a collision type and a mitigation directly, under two conditions: with the case's signals shown (99.3\% accuracy, $\kappa$=0.99) and with only the raw prompt (98.0\%, $\kappa$=0.98). This experiment is construct-validity evidence, directly comparable to the human-annotator numbers on the identical cases -- it does not feed Table~\ref{tab:governance-tradeoffs}. \texttt{llm\_detected\_edgecase\_routed.py} is the full-pipeline experiment behind Section~\ref{sec:llm-detected} and Table~\ref{tab:governance-tradeoffs}'s sixth column: the LLM classifies only a collision type (never a mitigation), which is then routed through the unmodified \texttt{INCOMPATIBLE\_ACTIONS} table, scored across all 10 statistical-rigor seeds. Because the benchmark generator draws prompts and signals from small fixed pools, the 12,600 (seed, case) rows reduce to only 803 distinct (prompt, signal-set) pairs; each was classified exactly once and the result cached and applied to every row sharing that pair, both to bound API cost (803 vs.\ 12,600 calls) and to avoid pseudo-replication -- treating repeated queries of an identical input as independent evidence would inflate the apparent sample size without adding information.

Checking why the raw-prompt-only condition scores as high as it does surfaced a distinct finding: the canonical 1,260-case benchmark, and separately the full 12,600-row pool across the 10 statistical-rigor seeds, are each built from only 35 distinct prompt strings (the most-repeated string appears 400 times across the pooled seeds), and every one of the 35 maps to exactly one collision type with zero counterexamples, checked exhaustively. Table~\ref{tab:governance-tradeoffs} is unaffected, since it scores from \texttt{signals} tags, not prompt text -- but it materially qualifies how the raw-prompt-only LLM result should be read (Limitations).

Every number in this appendix is reproducible via \texttt{experiments/rigid\_baseline\_action\_scoring.py}, \texttt{experiments/undetected\_penalty\_analysis.py}, \texttt{experiments/severity\_weight\_sensitivity.py}, \texttt{experiments/robustness\_matrix.py}, \texttt{experiments/energy\_proxy\_sensitivity.py}, \texttt{experiments/llm\_detector\_baseline.py}, \texttt{experiments/llm\_detected\_edgecase\_routed.py}, \texttt{experiments/template\_diversity\_analysis.py}, and \texttt{experiments/specification\_bug\_analysis.py} (which reproduces this benchmark's own generation order and RNG trajectory exactly, rather than an approximate re-simulation, to verify Section~\ref{sec:results}'s before/after detection rates), wired into \texttt{run\_all.sh} where their inputs don't require paid API calls or human annotation.

\section{Human Annotator Agreement}
\label{app:annotator-agreement}
To check whether an independent reader reaches the same judgment as the benchmark's own template-generated labels, two annotators independently labeled a blinded, stratified sample of 152 cases (22 per domain, capped at 20 for moderation) drawn from the held-out test split, seeing only each case's prompt and signal list -- not its domain or reference label -- plus a rubric sheet defining the seven collision types, their triggering obligations, and expected mitigations.

Annotator 1 matched the benchmark's ground-truth collision type on 88.8\% of cases (Cohen's $\kappa$=0.87 against the key); Annotator 2 matched on 100\% ($\kappa$=1.0). Inter-annotator agreement was 88.8\% ($\kappa$=0.87). One important caveat on how to read these numbers: the rubric sheet gives annotators the same per-obligation trigger-signal lists the registry itself uses (e.g.\ \texttt{care.escalate} triggers on \texttt{self\_harm, crisis, coercion, abuse\_disclosure}, listed verbatim on the rubric), and every case's shown signals are drawn from exactly the two obligations' trigger vocabularies by construction. A sufficiently careful annotator could therefore resolve most rows by table lookup against the rubric rather than independent semantic judgment of the raw scenario -- this measures whether a careful reader can correctly apply our disclosed labeling scheme, which is a real and useful property (it is evidence the scheme is legible and reproducible), but a narrower claim than blind independent judgment from the scenario alone.

Consistent with that caveat, Annotator 1's 17 disagreements with the key are not random: all 17 cluster in two systematic confusions rather than scattering across categories. Ten (of 17) confuse block-versus-escalate with privacy-versus-safeguarding in the crisis-support domain, and four confuse privacy-versus-safeguarding with memory-care-versus-memory-poisoning in public-benefits. Both confusions have a real cause in the registry, not annotator carelessness: \texttt{care.escalate} (block-versus-escalate's care-side obligation) and \texttt{safeguarding.preserve\_context} (privacy-versus-safeguarding's safeguarding-side obligation) share almost the same trigger vocabulary (\texttt{self\_harm}, \texttt{coercion}, \texttt{abuse\_disclosure} trigger both), so a case is only disambiguated by which \emph{other} signal is also present -- an easy distinction to miss under time pressure, and a genuine, disclosed source of representational overlap in the registry's design that happens to sit squarely in the care/crisis-support domain this paper's redistribution story depends on most. Full per-row disagreement data: \texttt{experiments/results/annotation\_agreement.json}.

\section{Additional Figures}
\label{app:additional-figures}
Figures~\ref{fig:cross-model} and~\ref{fig:ablation} below plot the results reported in Table~\ref{tab:cross-model} (Cross-model Consistency) and Table~\ref{tab:ablation} (Ablation Analysis), respectively.

\begin{figure*}[t]
\centering
\includegraphics[width=\textwidth]{figures/fig4_cross_model_comparison.pdf}
\caption{Cross-model comparison on 140 real API cases per provider (20/domain). Governance externality does not track latency: Qwen2.5-7B ties Claude Sonnet on $X_k$ despite roughly 5$\times$ the latency, since the energy proxy is derived from token/call counts rather than wall-clock time.}
\label{fig:cross-model}
\end{figure*}

\begin{figure}[h]
\centering
\includegraphics[width=\columnwidth]{figures/fig3_ablation_analysis.pdf}
\caption{Ablation analysis (single canonical seed, 1,260 cases). Removing the obligation registry or collision detection collapses mitigation accuracy to zero; removing runtime instrumentation has no effect, since mitigation selection depends only on trace signals, not instrumentation metadata.}
\label{fig:ablation}
\end{figure}

\makeatletter\global\let\isChecklistMainFile\@empty\makeatother
\input{ReproducibilityChecklist.tex}

\end{document}
